\begin{lemma}[Positive sector dimensions]%
    \label{lem:peps_positive_dimension_bounds}
    \lean{TNLean.PEPS.positive_dimension_bounds}
    \leanok
    For any finite family of positive integers $d_i$, the single-copy index
    space has cardinality $\sum_i d_i$, and
    \begin{align}
        \sum_i d_i\leq\sum_i d_i^2,
        \qquad
        \sum_i d_i<\sum_i d_i^2\ \Longleftrightarrow\ \exists i,\ d_i>1.
    \end{align}
\end{lemma}
\begin{proof}\leanok
    Each positive integer satisfies $d_i\leq d_i^2$, with equality exactly
    when $d_i=1$. Summing these inequalities gives the claims.
\end{proof}

\begin{theorem}[Matrix coordinates of the finite-group Fourier decomposition]%
    \label{thm:peps_regular_fourier_matrix}
    \lean{TNLean.PEPS.exists_minimalSemiRegular_leftRegular_matrix,
        TNLean.PEPS.card_eq_sum_sq_of_regular_intertwiner,
        TNLean.PEPS.bondDimensions_of_regular_intertwiner}
    \leanok
    \uses{thm:peps_regular_minimal_representation,
        lem:peps_positive_dimension_bounds}
    For every finite group $G$, there are positive dimensions $d_i$,
    pairwise inequivalent irreducible unitary matrix representations $D_i$,
    and an isometric matrix $Q$ such that
    \begin{align}
        U_{\mathrm{reg}}(g)
            =Q\left(\bigoplus_i D_i(g)\otimes I_{d_i}\right)Q^\dagger
        \qquad(g\in G).
    \end{align}
    The same matrices give a unitary semi-regular representation
    $U_{\min}=\bigoplus_iD_i$ with character multiplicity one in each sector.
    The dimensions satisfy $|G|=\sum_i d_i^2$ and
    $\sum_i d_i\leq |G|$, with strict inequality exactly when some $d_i>1$.
    The representation and coordinate matrix depend only on $G$.
    This is the matrix form of the Fourier decomposition used in
    \cite[Section~7]{Schuch2010PEPS}.
\end{theorem}

\begin{proof}\leanok
    Choose the orthonormal Fourier basis of the regular representation.
    Let $Q$ consist of its columns in the original group basis. Both bases
    are orthonormal, so $Q^\dagger Q=I$, and its conjugate transpose is the
    inverse change of coordinates. Applying the change-of-basis identity
    to each group matrix gives the displayed formula. The single-copy sum
    has the unitary, semi-regular, and multiplicity properties already proved.
    The identity group element also gives $QQ^\dagger=I$, so the two index
    spaces have the same dimension. Counting the repeated indices gives
    $|G|=\sum_i d_i^2$; the preceding numerical lemma gives the comparison.
\end{proof}
